\documentclass{article}
\usepackage{amsmath, amssymb, amsthm}
\usepackage{hyperref}
\usepackage{geometry}
\geometry{margin=1in}

\title{Erd\H{o}s Problem \#108}
\date{Last updated: 2025-08-31}
\author{Source: erdosproblems.com}

\begin{document}
\maketitle

\noindent\textbf{Status:} OPEN \quad \textbf{No prize}

\noindent\textbf{Tags:} graph theory, chromatic number, cycles

\medskip

\noindent\textbf{Problem Statement:}

For every $r\geq 4$ and $k\geq 2$ is there some finite $f(k,r)$ such that every graph of chromatic number $\geq f(k,r)$ contains a subgraph of girth $\geq r$ and chromatic number $\geq k$?




Conjectured by Erd\H{o}s and Hajnal. R\"{o}dl \cite{Ro77} has proved the $r=4$ case (see [923] ). The infinite version (whether every graph of infinite chromatic number contains a subgraph of infinite chromatic number whose girth is $&gt;k$) is also open. In \cite{Er79b} Erd\H{o}s also asks whether\[\lim_{k\to \infty}\frac{f(k,r+1)}{f(k,r)}=\infty.\]See also the entry in the graphs problem collection and [740] for the infinitary version.




References

[Er79b] Erd\H{o}s, Paul, Problems and results in graph theory and combinatorial analysis . Graph theory and related topics (Proc. Conf., Univ. Waterloo, Waterloo, Ont., 1977) (1979), 153-163.

[Ro77] R\"{o}dl, V., On the chromatic number of subgraphs of a given graph . Proc. Amer. Math. Soc. (1977), 370-371.

\noindent\textbf{Related OEIS sequences:} possible


\bigskip
\noindent\small{Source: \url{https://www.erdosproblems.com/108}}
\end{document}
